\section{Algoritmi bazë i një eksperimenti numerik}
Algoritmi kryesor i kapitullit nuk është një formulë e vetme, por cikli i plotë i punës. Ai duhet të përdoret në çdo notebook të kursit.

\begin{enumerate}
\item Formulohet pyetja fizike dhe përcaktohet madhësia e daljes.
\item Shënohen supozimet, njësitë, parametrat dhe kushtet fillestare.
\item Zgjidhet modeli minimal që mund t'i përgjigjet pyetjes.
\item Përcaktohen rrjeti, algoritmi, toleranca dhe kriteri i ndalimit.
\item Shkruhet kodi si funksione të testueshme.
\item Verifikohet me një rast analitik ose një ligj ruajtjeje.
\item Kryhet studimi i konvergjencës ose i numrit të mostrave.
\item Raportohen rezultati, gabimi numerik, pacaktueshmëria dhe kufijtë e modelit.
\end{enumerate}

\subsection{Një skelet Python për notebook-un}
\begin{lstlisting}[language=Python,caption={Skelet minimal për një eksperiment numerik.}]
from dataclasses import dataclass
import numpy as np
import matplotlib.pyplot as plt

@dataclass(frozen=True)
class Parametra:
    y0: float = 1.5       # m
    v0: float = 8.0       # m/s
    g: float = 9.81       # m/s^2
    t_fund: float = 1.2   # s
    dt: float = 0.02      # s

def modeli_analitik(t, p):
    """Pozicioni vertikal pa rezistence ajri."""
    return p.y0 + p.v0*t - 0.5*p.g*t**2

def eksperiment(p):
    t = np.arange(0.0, p.t_fund + p.dt/2, p.dt)
    y = modeli_analitik(t, p)
    assert np.isclose(y[0], p.y0)
    return t, y

p = Parametra()
t, y = eksperiment(p)
plt.plot(t, y)
plt.xlabel("Koha t [s]")
plt.ylabel("Lartesia y [m]")
plt.grid(alpha=0.25)
plt.show()
\end{lstlisting}

\subsection{Kontrolli i riprodhueshmërisë}
Për simulime stokastike përdoret një gjenerator i deklaruar, jo gjendja globale e fshehur.
\begin{lstlisting}[language=Python]
rng = np.random.default_rng(2026)
mostra = rng.normal(loc=0.0, scale=1.0, size=10_000)
\end{lstlisting}
Fara nuk e bën rezultatin ``më të rastit''; ajo identifikon realizimin. Në raport shënohen fara, numri i mostrave dhe versionet e bibliotekave.

\subsection{Terminologji dhe stil shkencor}
Në këtë tekst përdoren \emph{zbatim} për implementation, \emph{rrjetë} për grid, \emph{hap} për step, \emph{mbetje} për residual, \emph{vleftësim} për validation dhe \emph{riprodhueshmëri} për reproducibility. Kodi nuk duhet të fshehë njësitë fizike; ato shënohen në komente, tabela dhe boshte.

\begin{lidhje}
Notebook-u i studentit duhet të nisë nga ky skelet dhe të ruajë ndarjen ndërmjet parametrave, modelit, eksperimentit dhe paraqitjes grafike.
\end{lidhje}

Për organizimin terminologjik të modelit, algoritmit dhe zbatimit numerik janë konsultuar tekstet universitare shqip të analizës numerike \cite{hanelli,hoxha,datja}, krahas literaturës ndërkombëtare të fizikës llogaritëse \cite{newman,landau}.
